\documentclass[10pt]{article}
\usepackage[margin=0.72in]{geometry}
\usepackage{amsmath,amssymb}
\usepackage[hidelinks]{hyperref}
\setlength{\parindent}{0pt}
\setlength{\parskip}{0.35em}

\begin{document}
\begin{center}
{\LARGE KSGNS Universality and Paschke Dilations}\\[0.25em]
{\Large A Literature-Implied Answer}\\[0.55em]
{\large Status note for arXiv:1803.01911}\\[0.25em]
August 9, 2026
\end{center}

\begin{center}
\fbox{\parbox{0.90\textwidth}{\textbf{Status.}
Literature-implied answer, full in the von Neumann/self-dual-module setting
where the source defines Paschke dilations. This is not a new result of the
run.}}
\end{center}

\section*{Source question}

Westerbaan's thesis \cite{Westerbaan2019}, PDF page 57, paragraph 155 III,
states the KSGNS theorem for a completely positive map
\(\varphi\colon A\to\mathcal L_B(X)\), then asks whether KSGNS admits a
Stinespring-like universal property and whether it is a Paschke dilation.

The second clause is meaningful in the thesis's Paschke sense when the
algebras are von Neumann algebras, the maps are normal, and the Hilbert
modules are self-dual, so that their adjointable endomorphism algebras are von
Neumann algebras.

\section*{Identification}

Allen and Verdon \cite{AllenVerdon2024} supply exactly the required
von Neumann-module formulation. Their Proposition 1.2 realizes an arbitrary
von Neumann algebra as \(\operatorname{End}(X)\) for a generating
W*-module. Their Theorem 1.3 (PDF pp. 3--4, with proof beginning p. 17)
states that every normal completely positive map
\[
 f\colon \operatorname{End}(X)\longrightarrow\operatorname{End}(Y)
\]
has a bimodule Stinespring/KSGNS representation \((E,V)\). More decisively,
it states that a minimal representation is an \emph{initial object} in the
dagger category of representations. Consequently it is unique up to unitary
equivalence and admits a unique comparison isometry to every other
representation. This is the requested Stinespring-like universal property.

On PDF page 7, the same paper explicitly relates this construction to
Paschke dilations: for a minimal representation, the triple
\[
 \bigl(\operatorname{End}(X\mathbin{\bar\otimes}E),\rho,
       \operatorname{ad}_V\bigr)
\]
is a Paschke dilation. Thus both clauses of the source question have a
positive answer in the full setting in which ``Paschke dilation'' is defined.

\section*{Provenance and scope}

This is an agent-identified implication. Allen and Verdon cite the earlier
Paschke-dilation paper but do not state that they are answering the question
on page 57 of arXiv:1803.01911. Their result covers arbitrary von Neumann
algebras through generating W*-modules, rather than only type-I operator
algebras. For general C*-modules, minimal KSGNS uniqueness is standard, but
the thesis's normal-von-Neumann definition of a Paschke dilation does not
apply literally; no claim about a distinct C*-categorical analogue is made
here.

Bounded searches on August 9, 2026 used the exact source sentence, KSGNS,
``universal property'', ``Paschke dilation'', and the source arXiv id. They
found the Allen--Verdon theorem and no reason to regard the W*-case as open.

{\small
\begin{thebibliography}{9}
\bibitem{Westerbaan2019}
B.~E. Westerbaan,
\emph{Dagger and Dilation in the Category of Von Neumann Algebras},
arXiv:1803.01911v3, 2019.

\bibitem{AllenVerdon2024}
R.~P. Allen and D.~T. Verdon,
\emph{$\mathrm{CP}^{\infty}$ and beyond: 2-categorical dilation theory},
arXiv:2310.15776v3, 2024.
\end{thebibliography}
}

\end{document}
